\documentclass[10pt,a4paper]{amsart}
\usepackage[margin=19mm]{geometry}
\usepackage{amsmath,amssymb,mathtools}
\usepackage[scr=rsfs,cal=euler]{mathalfa}
\usepackage{xfrac}
\usepackage[unicode,hidelinks,bookmarks=false]{hyperref}
\newcommand{\ie}{i.e.\ }
\newcommand{\cf}{cf.\ }
\newcommand{\resp}{respectively\ }
\newcommand{\Subsection}{\subsection}
\providecommand{\nobreakdash}{\nobreak\hbox{-}}
\setlength{\emergencystretch}{2em}
\multlinegap0pt
\usepackage{amssymb}
\usepackage[shortlabels]{enumitem}
\usepackage{cancel}
\usepackage[normalem]{ulem}
\usepackage{mathtools}
\usepackage{xcolor}
\newtheorem{theorem}{Theorem}
\newcommand{\APQ}{\mathcal{H}_{P, Q}}
\title{Sharp Time-Decay Estimates for Fractional Heat Semigroups Associated with Polynomial Anharmonic Oscillators}
\author{Julio Delgado, Vishvesh Kumar, Shyam Swarup Mondal}
\date{Source: arXiv:2607.17580v1 (CC BY 4.0)}
\begin{document}
\maketitle

\noindent\emph{Source and licence.} Sharp Time-Decay Estimates for Fractional Heat Semigroups Associated with Polynomial Anharmonic Oscillators. Version arXiv:2607.17580v1 (CC BY 4.0), distributed by arXiv under the Creative Commons Attribution 4.0 International licence, http://creativecommons.org/licenses/by/4.0/, as recorded in the arXiv licence metadata for this exact version and shown on the abstract page https://arxiv.org/abs/2607.17580v1. Full licence notice: \texttt{licenses/arxiv-2607.17580-CC-BY-4.0.txt}.\par\smallskip

\noindent\textit{Editorial scope.} Counted unit: the article's theorem criticalglobal (source0.tex lines 2049--2057). The article's complete original proof of it is delivered with it (source0.tex lines 2059--2316, one \texttt{proof} environment of 258 lines). No mathematics is written, repaired or re-derived: the statement and the proof are the article's own lines, in the article's own notation, and no retained line is substituted or re-flowed.  Retained uncounted as the source-local context the proof needs: lines 354--404; lines 1745--1751.  Nothing was omitted from any retained span, so the package carries no omission record.  No retained line contains a citation, so the wrapper carries no bibliography and no bibliography entry.  No retained line contains a figure environment, an image-inclusion command or drawing code, so no image is delivered and the package declares no asset dependency.  Editorial formatting only, in addition: an AMS \texttt{amsart} preamble that restates the article's own declarations and macros used by the retained lines, namely the environments \texttt{theorem} on separate counters, together with the standard packages those lines need.  The retained environments print in this document as Theorem 1 in order of appearance, whereas the article prints the same items with the numbering of its own full document.  The complete native source of the article as distributed in the arXiv e-print for this exact version is bundled unchanged as \texttt{sources/205580/source0.tex}.
\begin{theorem}
\label{Lpmapping}
Let $s>0$. Then the following estimates hold.

\begin{itemize}
\item If $t>1$, then, for all $p,q\in[1,\infty]$,
\begin{equation}
\label{lpest}
    \Vert e^{-t\APQ^s}f\Vert_{L^q(\mathbb R^n)}
    \leq
    C e^{-t\lambda_0^s}
    \Vert f\Vert_{L^p(\mathbb R^n)}.
\end{equation}

\item If $0<t\leq1$, then
\begin{equation}
\label{lpest1}
    \Vert e^{-t\APQ^s}f\Vert_{L^q(\mathbb R^n)}
    \leq
    C t^{-\sigma_s}
    \Vert f\Vert_{L^p(\mathbb R^n)},
\end{equation}
where $\sigma_s$ is given as follows:
\begin{itemize}
\item[(i)] If $p,q\in(1,\infty)$, then
\[
    \sigma_s=
    \begin{cases}
    \displaystyle
    \frac{n}{2k s}
    \left(\frac1q-\frac1p\right),
    & q\leq p,\\[3mm]
    \displaystyle
    \frac{n}{2\ell s}
    \left(\frac1p-\frac1q\right),
    & p\leq q.
    \end{cases}
\]

\item[(ii)] If $p=1$ and $q=\infty$, then $ \sigma_s=\frac{n}{2\ell s}.$

\item[(iii)] If $p=1$ and $2\leq q<\infty$, then $ \sigma_s=
    \frac{n}{2\ell s}
    \left(1-\frac1q\right).$

\item[(iv)] If $1<p<\infty$ and $q=1$, then $ \sigma_s=
    \frac{n}{2k s}
    \left(1-\frac1p\right).$
\end{itemize}
\end{itemize}
\end{theorem}

\begin{equation}
\label{eq6.1}
    \begin{cases}
        \partial_t u+\APQ^s u=|u|^{\beta-1}u,\\
        u(0,x)=u_0(x),
    \end{cases}
\end{equation}

\begin{theorem}
\label{criticalglobal}
Let $s>0$, $\beta>1$, and assume that
\[
    p_c^s:=\frac{n(\beta-1)}{2\ell s}>1.
\]
Let $u_0\in L^{p_c^s}(\mathbb R^n)$. If $\Vert u_0\Vert_{L^{p_c^s}(\mathbb R^n)}$
is sufficiently small, then the corresponding solution of \eqref{eq6.1} is global, that is, $T_{\max}=\infty.$
\end{theorem}

\begin{proof}
Choose $r$ such that $p_c^s<r<\beta p_c^s.$
Equivalently, $ \frac{1}{\beta p_c^s}
<\frac1r<
    \frac1{p_c^s}.$
Since $ p_c^s=\frac{n(\beta-1)}{2\ell s},$
this is the same as $ \frac{2\ell s}{n\beta(\beta-1)} <
    \frac{1}{r}   <
    \frac{2\ell s}{n(\beta-1)}.$
Define $\rho :=
    \frac{1}{\beta-1} -
    \frac{n}{2r\ell s}.$
Then $\rho>0$. Moreover, we see that 
\begin{equation}
\label{lessone}
    \rho +1-
    \frac{n(\beta-1)}{2r\ell s}-
    \rho\beta =0.
\end{equation}
We also note that
\[
    \rho\beta<1,
    \qquad
    \frac{n(\beta-1)}{2r\ell s}<1.
\]

Let
\[
    Y
    :=
    \left\{
    u:(0,\infty)\to L^r(\mathbb R^n):
    \Vert u\Vert_Y<\infty
    \right\},
\]
where $\Vert u\Vert_Y
    :=
    \sup_{t>0}
    t^\rho
    \Vert u(t)\Vert_{L^r}.$
    
For $M>0$, set $Y_M
    :=
    \left\{
    u\in Y:
    \Vert u\Vert_Y\leq M
    \right\},$
equipped with the metric $$d(u,v)
    :=
    \sup_{t>0}
    t^\rho
    \Vert u(t)-v(t)\Vert_{L^r}.$$
Then $(Y_M,d)$ is a complete metric space.

Define the fixed-point map
\[
    \mathfrak J_{u_0}(u)(t)
    :=
    e^{-t\APQ^s}u_0
    +
    \int_0^t
    e^{-(t-\tau)\APQ^s}
    \big(|u(\tau)|^{\beta-1}u(\tau)\big)
    \,d\tau.
\]
Assume first that
\begin{equation}
\label{6.11}
    \sup_{t>0}
    t^\rho
    \Vert e^{-t\APQ^s}u_0\Vert_{L^r}
    \leq
    \delta.
\end{equation}
We shall choose $\delta>0$ sufficiently small later.

Let $u,v\in Y_M$. Using Theorem \ref{Lpmapping} with exponents $ \left(\frac r\beta,r\right),$
we obtain
\[
\begin{aligned}
&\Vert
e^{-(t-\tau)\APQ^s}
\big(
|u(\tau)|^{\beta-1}u(\tau)
-
|v(\tau)|^{\beta-1}v(\tau)
\big)
\Vert_{L^r}
\\
&\qquad
\leq
C
(t-\tau)^{-\frac{n(\beta-1)}{2r\ell s}}
\Vert
|u(\tau)|^{\beta-1}u(\tau)
-
|v(\tau)|^{\beta-1}v(\tau)
\Vert_{L^{r/\beta}}.
\end{aligned}
\]
By H\"older's inequality and the pointwise Lipschitz bound for the nonlinearity,
\[
\begin{aligned}
\Vert
|u|^{\beta-1}u
-
|v|^{\beta-1}v
\Vert_{L^{r/\beta}}
\leq
C_\beta
\left(
\Vert u\Vert_{L^r}^{\beta-1}
+
\Vert v\Vert_{L^r}^{\beta-1}
\right)
\Vert u-v\Vert_{L^r}.
\end{aligned}
\]
Since $u,v\in Y_M$, $\Vert u(\tau)\Vert_{L^r} +
    \Vert v(\tau)\Vert_{L^r}
    \leq C M\tau^{-\rho},$
and $\Vert u(\tau)-v(\tau)\Vert_{L^r}
    \leq
    d(u,v)\tau^{-\rho},$ we have

\[
\begin{aligned}
\Vert
e^{-(t-\tau)\APQ^s}
\big(
|u(\tau)|^{\beta-1}u(\tau)
-
|v(\tau)|^{\beta-1}v(\tau)
\big)
\Vert_{L^r}
\leq
C
M^{\beta-1}
d(u,v)
(t-\tau)^{-\frac{n(\beta-1)}{2r\ell s}}
\tau^{-\rho\beta}.
\end{aligned}
\]
Thus,
\[
\begin{aligned}
t^\rho
\Vert
\mathfrak J_{u_0}(u)(t)
-
\mathfrak J_{u_0}(v)(t)
\Vert_{L^r}
\leq
C
t^\rho
M^{\beta-1}
d(u,v)
\int_0^t
(t-\tau)^{-\frac{n(\beta-1)}{2r\ell s}}
\tau^{-\rho\beta}
\,d\tau.
\end{aligned}
\]
Using the change of variables $\tau=t\theta$, we get
\[
\begin{aligned}
t^\rho
\int_0^t
(t-\tau)^{-\frac{n(\beta-1)}{2r\ell s}}
\tau^{-\rho\beta}
\,d\tau
=
t^{
\rho+1
-\frac{n(\beta-1)}{2r\ell s}
-\rho\beta
}
\int_0^1
(1-\theta)^{-\frac{n(\beta-1)}{2r\ell s}}
\theta^{-\rho\beta}
\,d\theta.
\end{aligned}
\]
The integral is finite because $ \rho\beta<1,\,
    \frac{n(\beta-1)}{2r\ell s}<1,$
and the power of $t$ is zero by \eqref{lessone}. Hence there exists $L>0$ such that
\begin{equation}
\label{criticalcontraction}
    d\big(\mathfrak J_{u_0}(u),\mathfrak J_{u_0}(v)\big)
    \leq
    L M^{\beta-1}d(u,v).
\end{equation}
Similarly, taking $v=0$, we obtain
$\Vert \mathfrak J_{u_0}(u)\Vert_Y
    \leq \delta + L M^\beta.$

Choose $M>0$ so small that $L M^{\beta-1}\leq \frac12.$
Then choose $\delta>0$ so small that
\[
    \delta+L M^\beta\leq M.
\]
With these choices, $\mathfrak J_{u_0}$ maps $Y_M$ into itself and is a strict contraction on $Y_M$. Therefore, by Banach's fixed point theorem, there exists a unique global fixed point $u\in Y_M$.

It remains to verify that \eqref{6.11} follows from smallness of $u_0$ in $L^{p_c^s}$. By Theorem \ref{Lpmapping}, for $0<t\leq1$,
\[
    \Vert e^{-t\APQ^s}u_0\Vert_{L^r}
    \leq
    C
    t^{-\frac{n}{2\ell s}
    \left(
    \frac1{p_c^s}-\frac1r
    \right)}
    \Vert u_0\Vert_{L^{p_c^s}}.
\]
Since $\frac{n}{2\ell s}
    \left(
    \frac1{p_c^s}-\frac1r
    \right)  =
    \rho,$
we obtain
\[
    t^\rho
    \Vert e^{-t\APQ^s}u_0\Vert_{L^r}
    \leq
    C
    \Vert u_0\Vert_{L^{p_c^s}},
    \qquad 0<t\leq1.
\]
For $t>1$, the large-time estimate in Theorem \ref{Lpmapping} gives
\[
    \Vert e^{-t\APQ^s}u_0\Vert_{L^r}
    \leq
    C e^{-t\lambda_0^s}
    \Vert u_0\Vert_{L^{p_c^s}},
\]
and hence
\[
    t^\rho
    \Vert e^{-t\APQ^s}u_0\Vert_{L^r}
    \leq
    C
    t^\rho e^{-t\lambda_0^s}
    \Vert u_0\Vert_{L^{p_c^s}}
    \leq
    C
    \Vert u_0\Vert_{L^{p_c^s}}.
\]
Therefore,
\[
    \sup_{t>0}
    t^\rho
    \Vert e^{-t\APQ^s}u_0\Vert_{L^r}
    \leq
    C
    \Vert u_0\Vert_{L^{p_c^s}}.
\]
Thus \eqref{6.11} holds whenever $\Vert u_0\Vert_{L^{p_c^s}}$ is sufficiently small. Consequently, the solution is global. This completes the proof.
\end{proof}

\end{document}
